\section{Lower jets and the determinant contradiction}
\label{sec:jets}

The comparison maxima of Section~\ref{sc:section} relate a center
endpoint to finitely many active outer endpoints. We seek determinant
bounds for lower tests of the original potential at these endpoints.
The main obstacle is that the center selected by a maximum need not
belong to the full-measure set where the original transport Jacobian
equation holds. Even though the short-time minimizing map is Lipschitz,
it could send a set of positive volume into that exceptional set.

We first transfer lower jets from the prefix regularization to its
active outer endpoints. The original contact-volume bound applies to
these lower tests and bounds the selected outer determinant. That bound
controls the regularized center Hessian, which in turn makes the
short-time minimizing map nonsingular at the sampled points. Only then
can we use the original Jacobian equation at the centers. The resulting
matrix inequalities give a determinant gain proportional to the outer
curvature parameter~$K$, contradicting the geometric comparison from
Section~\ref{sec:spectral}.

Throughout the argument the manifold and density bounds are fixed.
In the local calculations, $F,G$ and their indexed versions denote
regularized potentials; the scalar oscillation function $F$ from
Section~\ref{sc:section} returns in Proposition~\ref{jet:contradiction}.

Throughout this section, a \emph{lower Taylor jet} $(\zeta,B)$ of a function $f$ at the
origin means
\[
 f(h)\ge f(0)+\zeta\cdot h+\tfrac12 B(h,h)+o(|h|^2).
\]
Replacing $B$ by $B-\eps\Id$, for any $\eps>0$, gives a genuine quadratic
lower test on a sufficiently small neighborhood. All the tests to which
we apply a determinant inequality will have a strictly positive relative
Hessian, so this replacement and the subsequent limit $\eps\downarrow0$
are harmless.

\subsection{The original contact measure and lower tests}

\begin{lemma}[Contact-volume comparison]\label{jet:contact-volume}
Let $u,v$ be the dual potentials for an admissible pair of densities,
and let
\[
 \mathcal C=\{(x,y):u(x)+c(x,y)+v(y)=0\}.
\]
For a Borel set $A$ on either side, its contact image satisfies
\begin{equation}\label{jet:contact-volume-eq}
 \frac{\lambda}{\Lambda}\vol(A)
 \le \vol(\mathcal C(A))
 \le \frac{\Lambda}{\lambda}\vol(A).
\end{equation}
Projections are interpreted in the completed volume measures.
\end{lemma}

\begin{proof}
The contact relation is closed. Its projections of Borel sets are
analytic, hence measurable for the completed measures. The optimal
coupling has full mass on $\mathcal C$. At every differentiability point
of either potential its contact endpoint is unique: all minimizing
first-variation logs must equal its gradient, and their exponential
endpoints therefore coincide. Since both marginals are absolutely
continuous, almost every contact pair lies in the sets where both
potentials are differentiable. The two contact maps on these sets are
inverse to each other.

Write $\mu=\rho_0\vol$ and $\nu=\rho_1\vol$. Consequently
$\nu(\mathcal C(A))=\mu(A)$ for sets on the first side, and
$\mu(\mathcal C(A))=\nu(A)$ for sets on the second side: any additional
contact image outside the almost-everywhere map image consists, up to a
null set, of points where the opposite potential is not differentiable.
The density bounds give \eqref{jet:contact-volume-eq}.
\end{proof}

\begin{lemma}[Upper determinant bound for a lower test]
\label{jet:viscosity}
Let $v$ be one of the original dual potentials. Suppose a $C^2$
function $\phi$ touches $v$ from below at $y$, its gradient
$r=\grad\phi(y)$ belongs to $I(y)$, and
\[
 W=\Hess\phi(y)+A_y(r)>0.
\]
Then
\begin{equation}\label{jet:viscosity-eq}
 \sigma_y(r)\det W\le \Lambda/\lambda.
\end{equation}
The same conclusion holds for a lower Taylor jet with these
properties, by approximation of its quadratic part.
\end{lemma}

\begin{proof}
It suffices first to prove the result after decreasing the Hessian by a
small positive multiple of the metric. Thus we may arrange strict
touching, with
$v-\phi\ge\eps d(\cdot,y)^2$ in a small coordinate ball. The map
\[
 S(z)=\exp_z\grad\phi(z)
\]
is a local diffeomorphism near $y$. Indeed its defining equation is
$D_zc(z,S(z))+\grad\phi(z)=0$, and differentiation gives
\begin{equation}\label{jet:smooth-jacobian}
 |\det DS(z)|
 =\sigma_z(\grad\phi(z))
   \det\{\Hess\phi(z)+A_z(\grad\phi(z))\}.
\end{equation}
The relative Hessian is positive throughout a sufficiently small ball,
and the logs remain before cut.

For $h>0$, let $\Omega_h$ be the connected component containing $y$ of
$\{\phi+h>v\}$ inside that ball. It is relatively compact and shrinks to
$y$. If $z\in\Omega_h$ and $x=S(z)$, the cost mountain
\[
 m_x(w)=\phi(z)+c(z,x)-c(w,x)
\]
is tangent to $\phi$ at $z$. On a fixed smaller coordinate ball,
$\phi-m_x$ has positive Hessian, uniformly for $x$ close to $S(y)$;
thus $\phi\ge m_x$ there. On $\partial\Omega_h$ we have
$v-m_x\ge h$, whereas at $z$ we have $v(z)-m_x(z)<h$.
The minimum of $v-m_x$ on $\overline{\Omega_h}$ is therefore attained
at an interior point $z'$. Thus $\grad m_x(z')\in\partial v(z')$.
By Theorem~\ref{geo:support}, the corresponding local contact is
a global cost contact. Its endpoint is $x$. This proves
\[
 S(\Omega_h)\subset\mathcal C(\Omega_h).
\]
Use the area formula for the diffeomorphism $S$
\cite[Chapter~2, Theorem~3.3]{Simon2018}, followed by
Lemma~\ref{jet:contact-volume}, to obtain
\[
 \inf_{\Omega_h}|\det DS|\,\vol(\Omega_h)
 \le \vol(S(\Omega_h))
 \le(\Lambda/\lambda)\vol(\Omega_h).
\]
Let $h\downarrow0$ and then restore the original Hessian.
\end{proof}

Only the original correspondence occurs in these two Lemmas.
No density bound is claimed for a modified potential or its envelope.
For later use, the almost-everywhere counterpart of
\eqref{jet:smooth-jacobian} is the following consequence of
\cite[Proposition~4.1 and Theorem~4.2]{CEMS2001}. There is a
volume-null Borel set $N\subset M$ such that, at every $y\notin N$,
$v$ has an Alexandrov Hessian, its true log $r=\grad v(y)$ is before
cut, and
\begin{equation}\label{jet:ae-jacobian}
 W_v(y):=\Hess v(y)+A_y(r)>0,\qquad
 \frac{\lambda}{\Lambda}
 \le \sigma_y(r)\det W_v(y)
 \le\frac{\Lambda}{\lambda}.
\end{equation}
We enlarge $N$ to include the exceptional sets in the contact-map
differential of \cite[Proposition~4.1(b)]{CEMS2001}. The Jacobian
equation in \cite[Theorem~4.2]{CEMS2001} assumes only absolute
continuity with $L^1$ densities. The lower density bound makes its
full-source-measure set a full-volume set; differentiability of a
density is not required.

\subsection{An exact endpoint-jet calculation}

Fix a continuous datum $w_o$ and put
\[
 g(x)=\max_y\{-c(x,y)-w_o(y)\},\qquad F=Q_tg,\quad 0<t<1.
\]
By Proposition~\ref{geo:intermediate}, $F$ is locally $C^{1,1}$ and its
minimizing pole $x(z)$ is unique and locally Lipschitz. The following
calculation concerns the first minimum alone.

\begin{lemma}[The pole matrix and endpoint lower jets]
\label{jet:transfer}
Suppose $z_0$ is a point where $F$ has a second-order expansion and
$x(z)$ is differentiable. Write $x=x(z_0)$,
$z_0=\exp_x(tp)$, and choose active logs
\[
 p=\sum_{j=1}^Nm_jp_j,\qquad m_j>0,\quad
 \sum_jm_j=1,\qquad N\le n+1,
\]
where $y_j=\exp_xp_j$ and
$g(x)+c(x,y_j)+w_o(y_j)=0$. Set
$E=\Span\{p_j-p:1\le j\le N\}$. In the fixed-pole coordinates
$z(q)=\exp_x(tq)$, define
\begin{equation}\label{jet:L-definition}
 L=D_q^2\{c_t(x,z(q))-F(z(q))\}\big|_{q=p}.
\end{equation}
Then
\begin{equation}\label{jet:L-properties}
 L\ge0,\qquad LE=0.
\end{equation}
For every $j$, the function
\[
 \Psi_j(d)=w_o(\exp_x(p_j+d))+\tfrac12|p_j+d|^2
\]
has at $d=0$ the lower Taylor jet $(0,m_jL)$, after subtracting
its value at zero. If $p_j$ is nonconjugate, this is an endpoint
lower test in the smooth chosen cost branch.

More precisely, for $t<s<1$ put
\[
 Q_s=\frac1tA_x(tp)-\frac1s\sum_jm_j A_x(sp_j).
\]
Then $Q_s\ge0$ on $E^\perp$, and
$m_j(L+LQ_sL)$ is a lower Taylor Hessian for $\Psi_j$.
\end{lemma}

\begin{proof}
Use normal coordinates $x'=\exp_xa$ and set
\[
 B(a,q)=c_t(\exp_xa,z(q)),\quad
 f(q)=F(z(q)),\quad a(q)=\log_xx(z(q)).
\]
At the reference point $a(p)=0$. First variation gives the exact
identities
\[
 B(0,q)=\tfrac t2|q|^2,\qquad B_a(0,q)=-q.
\]
Thus $B_{qq}=t\Id$, $B_{qa}=-\Id$ at $(0,p)$.
The envelope identity $Df(q)=B_q(a(q),q)$ and differentiability of the
pole imply
\begin{equation}\label{jet:pole-identity}
 D^2f=t\Id-R,\qquad R=D_qa(p)=L.
\end{equation}
In particular $R$ is symmetric. Since
$f(q)\le g(x)+B(0,q)$ with equality at $p$, we have $L\ge0$.

Choose a common $s\in(t,1)$. For each $j$ define
\[
 y_j(d)=\exp_x(p_j+d),\qquad
 z_j(d)=\exp_x(s(p_j+d)),
\]
and the split upper cost
\begin{equation}\label{jet:split-cost}
 C_j(a,d)=
 \frac{c(\exp_xa,z_j(d))}{s}
 +\frac{c(z_j(d),y_j(d))}{1-s}.
\end{equation}
Both cost pieces are smooth near $(0,0)$ because they are proper
subsegments of a minimizing ray. Energy splitting gives
$c(\exp_xa,y_j(d))\le C_j(a,d)$. At $a=0$, the chosen geodesic
remains minimizing on both pieces for small $d$, and consequently
\begin{equation}\label{jet:split-derivatives}
 C_j(0,d)=\tfrac12|p_j+d|^2,\quad
 (C_j)_a(0,d)=-(p_j+d),\quad
 (C_j)_{ad}=-\Id,\quad
 (C_j)_{aa}=\frac1sA_x(sp_j).
\end{equation}
The use of the fixed reference pole in $z_j(d)$ is what makes the
mixed derivative exactly $-\Id$.

The nonnegative function
\[
 g(x(z(q)))+C_j(a(q),0)+w_o(y_j)
\]
vanishes at $p$. It is differentiable there, since
$g(x(z(q)))=f(q)-B(a(q),q)$, and its derivative equals
$(p-p_j)^*R$. This derivative is zero. Symmetry in
\eqref{jet:pole-identity} proves $LE=0$.

On $E^\perp$, MTW concavity on the convex hull of the $sp_j$, followed
by divided-action monotonicity, gives
\[
 Q_s\ge \frac1tA_x(tp)-\frac1s A_x(sp)>0.
\]
All points used in this inequality are before cut by
Theorem~\ref{geo:support}; equivalently one may invoke
Lemma~\ref{geo:slack}. Notice that unsplit endpoints may still be
conjugate.

Fix $i$ and write $m=m_i$. Given $d$, choose any linear map
$U:T_xM\to E$, and set $k=m(d+Ud)$ and $a=a(p+k)$.
In the lower bound defining $g(\exp_xa)$, move endpoint $i$ to
$y_i(d)$, keep all other endpoints fixed, replace the costs by
\eqref{jet:split-cost}, and average with the old weights. Write
$\widehat\Psi_i(d)=\Psi_i(d)-\Psi_i(0)$. The resulting inequality is
\[
 f(p+k)\ge
 B(a,p+k)-\sum_jm_j\{C_j(a,d_j)+w_o(y_j(d_j))\},
\]
where $d_i=d$ and $d_j=0$ for $j\ne i$. Subtract the equality at
$(a,k,d)=(0,0,0)$ and use
\eqref{jet:pole-identity}--\eqref{jet:split-derivatives}. The linear
terms cancel: the pole terms cancel because $\sum m_jp_j=p$, and
the prefix terms cancel because $Df(p)=tp$. The remaining inequality is
\begin{equation}\label{jet:quadratic-transfer}
 m\widehat\Psi_i(d)\ge
 \tfrac12L(k,k)+\tfrac12Q_s(a,a)
 +a\cdot(md-k)+o(|d|^2).
\end{equation}
Here $a=Lk+o(|d|)$. Since $md-k=-mUd$ and
$\operatorname{range}L\subset E^\perp$, the mixed term is
$o(|d|^2)$. Hence the lower Hessian is
\[
 m(\Id+U)^*(L+LQ_sL)(\Id+U).
\]
In fact this equals $m(L+LQ_sL)$ because $LU=0$. The latter
dominates $mL$, proving the asserted jet. The linear term is exactly
zero independently of any second function or maximum condition.

The map $d\mapsto y_i(d)$ need not be invertible for the preceding
pullback calculation. When $p_i$ is nonconjugate it is a local
diffeomorphism, and $\tfrac12|p_i+d|^2$ is the smooth chosen
upper branch of $c(x,y_i(d))$. This converts the pullback jet into
an endpoint jet.
\end{proof}

Let $G$ be another locally $C^{1,1}$ function, and define
\begin{equation}\label{jet:H-definition}
 H=D_q^2\{c_t(x,z(q))+G(z(q))\}\big|_{q=p}.
\end{equation}
The exact identity
\begin{equation}\label{jet:H-minus-L}
 H-L=D_q^2\{(F+G)(z(q))\}\big|_{q=p}
\end{equation}
will distinguish the first-minimum calculation from the maximum
argument. At a twice differentiable maximum, it gives $H\le L$ and
$H|_E\le0$. At a Jensen point with gradient and upper Hessian errors
tending to zero in a fixed smooth chart, it gives
\begin{equation}\label{jet:Jensen-error}
 H\le L+\eta\Id,\qquad \eta\longrightarrow0.
\end{equation}
The coordinate-change term uses the gradient error and also tends to
zero. Lemma~\ref{jet:transfer} is exact at such a point. For every
linear $U$ with image in $E$, it therefore supplies a lower jet
$H_i=m_iL$ satisfying
\[
 H_i\ge m_i(\Id+U)^*H(\Id+U)
       -m_i\eta(\Id+U)^*(\Id+U).
\]
In particular no approximate stationarity term occurs in the
linear coefficient of the endpoint test.

\subsection{The center semibound without a density equation}

We use the scales and templates of the preceding section. To distinguish
a scalar argument from a point of the manifold, write
\[
 \varphi_c(h)=h+bB_c((h-a)/D),\quad
 \varphi_o(h)=h+bB_o((h-a)/D),\quad
 w_c=\varphi_c\circ v,\quad w_o=\varphi_o\circ v.
\]
For each fixed outer template the scales are sufficiently small that
both scalar maps have derivative between $1/2$ and $2$. On the fixed
compact interval of possible center levels, $B_c'$ is negative and
bounded away from zero. Its derivative bounds are fixed independently
of the outer curvature parameter $K$.

\begin{lemma}[A lower semibound before the center Jacobian]
\label{jet:semibound}
Fix one set of scales, one original potential $v$, and the two
templates. Let $t_j\uparrow1$ and $b_j\in[b_-,b_+]$. Put
\[
 F_j=Q_{t_j}g_{b_j},\qquad G_j=Q_{1-t_j}w_{c,b_j}.
\]
Let $z_j$ be simultaneous differentiability points for the objects in
Lemma~\ref{jet:transfer}, with second-order expansions of $G_j$, whose center levels lie in the
fixed compact interval above. Suppose
\[
 |\grad(F_j+G_j)(z_j)|\longrightarrow0.
\]
Let $x_j$ be the prefix pole and
$z_j=\exp_{x_j}(t_jp_j)$. The matrix $H_j$ of
\eqref{jet:H-definition} satisfies
\begin{equation}\label{jet:semibound-eq}
 H_j\ge c\frac{b_j}{D}\Id
       -C\frac{b_j}{D^2}p_jp_j^*-\eps_j\Id,
 \qquad \eps_j\longrightarrow0.
\end{equation}
The constants $c,C$ depend on the geometry and the fixed center
template, but not on $K$. The errors tend to zero with the prefix and
gradient errors at the fixed scales. No upper bound on $H_j$, and no
Alexandrov Hessian of $v$ at the center, is assumed.
\end{lemma}

\begin{proof}
Write $\tau_j=1-t_j$, and let $y_j$ be the unique final minimizer.
Set
\[
 r_j=\tau_j^{-1}\log_{y_j}z_j,\qquad
 l_j=\varphi_c'(v(y_j))<1.
\]
The short minimum gives a smooth local lower test for $w_c$ with
gradient $r_j$. Inverting $\varphi_c$ gives a local lower test for
$v$ with gradient $r_j/l_j$. Theorem~\ref{geo:support} therefore
implies that
\[
 \zeta_j=\exp_{y_j}(r_j/l_j)
\]
is a true contact endpoint of $v$, joined by the indicated minimizing
ray. In particular $\widehat x_j=\exp_{y_j}r_j$ is strictly before
cut along that ray.

For the moment omit subscripts. Put
$\theta=(1+l^{-1})/2$ and $\xi=\exp_y(\theta r)$.
Split the true support at $\xi$, using the fraction $l\theta$ of its
unit duration. The resulting upper cost support is
\[
 C(h)=\frac{c(h,\xi)}{l\theta}
       +\frac{c(\xi,\zeta)}{1-l\theta}.
\]
It agrees with $c(h,\zeta)$ at $h=y$, and $v(h)\ge
v(y)+c(y,\zeta)-C(h)$. Thus
\[
 \psi(h)=\varphi_c\{v(y)+c(y,\zeta)-C(h)\}
\]
is a smooth local lower support for $w_c$, with
\begin{equation}\label{jet:center-support-jet}
 \grad\psi(y)=r,\qquad
 \Hess\psi(y)
 =-\theta^{-1}A_y(\theta r)
   +\frac{\varphi_c''(v(y))}{l^2}rr^*.
\end{equation}
The original true endpoint $\zeta$ may be conjugate. The point $\xi$
is halfway into the available extension and is before cut, which is
all that this construction uses.

At the fixed scales, $1-l$ has a positive lower bound and the
parameters $(y,r,l,v(y))$ range over a compact set. The first cost
piece in $C$ is uniformly smooth on a fixed neighborhood of $y$;
the second piece is constant in $h$. Consequently the supports
$\psi$ have uniform smooth bounds and a uniform supporting radius.
For sufficiently small $\tau$, the local infimum $Q_\tau\psi$ is
smooth near $z$, has minimizer $y$ at $z$, and supports
$Q_\tau w_c$ from below there. To justify the last assertion, all
minimizers of the latter lie within $O(\tau)$ of $z$, by the
Lipschitz bound on $w_c$, and so lie in the supporting neighborhood.
For $\psi$ the local Hessian of the minimizing expression is
$\tau^{-1}\Id+O(1)>0$. Equality and stationary first variation hold
at $y$. The implicit-function theorem and the Schur formula for a
minimum give, after parallel identification,
\begin{equation}\label{jet:smooth-short-limit}
 \Hess Q_\tau\psi(z)=\Hess\psi(y)+O(\tau).
\end{equation}
All constants in this statement are uniform for the fixed-scale
family just described.

Let $\widehat p=\log_{\widehat x}y$ and
$\widehat L=(d\exp_{\widehat x})_{\widehat p}$.
Lemma~\ref{spec:pullback}, applied to the minimizing ray through
duration $1/l$, gives
\[
 \widehat L^*
 \{A_y(r)-\theta^{-1}A_y(\theta r)\}\widehat L
 \ge c(1-l)\Id.
\]
Gauss's Lemma gives $\widehat L^*r=-\widehat p$. Combining this
with \eqref{jet:center-support-jet} yields
\begin{equation}\label{jet:center-support-pullback}
 \widehat L^*\{\Hess\psi(y)+A_y(r)\}\widehat L
 \ge c(1-l)\Id
       +\frac{\varphi_c''(v(y))}{l^2}
          \widehat p\widehat p^*.
\end{equation}
Here $1-l\ge c b/D$ and
$|\varphi_c''(v(y))|/l^2\le Cb/D^2$.

We finally compare the reference pole with $\widehat x_j$.
The envelope identities give
\[
 x_j=\exp_{z_j}(-t_j\grad F_j(z_j)),\qquad
 \widehat x_j=\exp_{z_j}(t_j\grad G_j(z_j)).
\]
They are at distance tending to zero. The pairs
$(\widehat x_j,y_j)$ have a uniform fixed-scale extension, so their
cost charts and inverse exponential charts are uniformly smooth.
Omit subscripts again and put $U=Q_\tau\psi$, $Z(q)=\exp_x(tq)$,
$M=D_qZ(p)$, and $R=\nabla_q^2Z(p)$.
Since $G-U$ has a local minimum at $z$, its first derivative
vanishes and its second-order form is positive semidefinite.
Consequently
\[
 H\ge H_\psi
 :=t\Id+M^*\Hess U(z)M+\langle\grad U(z),R\rangle.
\]
Only the smooth support $U$ occurs on the right. Its gradient
and Hessian are uniformly bounded at the fixed scales, and
\eqref{jet:smooth-short-limit} identifies its Hessian with that
of $\psi$, up to $O(\tau)$ after parallel transport.

Write $\delta=|\grad(F+G)(z)|$ and identify nearby pole
tangent spaces by parallel transport. The explicit pole formulas
give $d(x,\widehat x)=O(\delta)$.
Moreover $p=t^{-1}\log_xz$ and
$\widehat p=t^{-1}\log_{\widehat x}z$ differ by
$O_{\mathrm{fixed}}(\delta)$, since the latter ray has a
uniform fixed-scale extension. Thus replacing $x,p$ by
$\widehat x,\widehat p$ in $M,R$ changes their entries by
$O_{\mathrm{fixed}}(\delta)$. Changing $t$ to $1$ then
changes them by $O(\tau)$. Every resulting error in $H_\psi$
multiplies a bounded derivative of $U$.

For the full endpoint map $Y(q)=\exp_{\widehat x}q$, let
$B=DY(\widehat p)$. The identity
$c(\widehat x,Y(q))=|q|^2/2$ on the chosen branch gives
\[
 \Id=B^*A_y(r)B-\langle r,\nabla_q^2Y(\widehat p)\rangle.
\]
Since $\grad\psi(y)=r$, the preceding expansions yield
\[
 H\ge H_\psi
 =B^*\{\Hess\psi(y)+A_y(r)\}B
       +O_{\mathrm{fixed}}(\tau+\delta).
\]
Combining this with \eqref{jet:center-support-pullback}
proves \eqref{jet:semibound-eq}. No coordinate error
multiplies the unknown Hessian of $G$.
\end{proof}

The true and modified endpoints occur in opposite orders along the
center and outer reverse rays. In the center calculation, $l<1$ makes $\widehat x=\exp_y r$ a proper intermediate
point on the true ray to $\zeta=\exp_y(r/l)$. At an active outer
endpoint $y_i$, a slope $l_i=\varphi_o'(v(y_i))>1$ instead makes the true endpoint
$\zeta_i=\exp_{y_i}(r_i/l_i)$ an intermediate point on the modified
ray to $x=\exp_{y_i}r_i$, where $r_i$ is its minimizing reverse log.
Figure~\ref{fig:reverse-rays} displays this distinction. We next use it
to exclude conjugacy at the full modified outer endpoint.

\begin{figure}[htbp]
\centering
\setlength{\unitlength}{1mm}
\begin{picture}(128,57)
 \put(0,50){\makebox(128,5)[l]{Center modification: $l<1$}}
 \put(9,38){\vector(1,0){110}}
 \put(13,38){\circle*{1.2}}
 \put(72,38){\circle*{1.2}}
 \put(107,38){\circle*{1.2}}
 \put(8,40){\makebox(10,5){$y$}}
 \put(67,40){\makebox(10,5){$\widehat x$}}
 \put(102,40){\makebox(10,5){$\zeta$}}
 \put(8,31){\makebox(10,5){$0$}}
 \put(67,31){\makebox(10,5){$1$}}
 \put(97,31){\makebox(20,5){$1/l>1$}}
 \put(0,22){\makebox(128,5)[l]{Selected outer modification: $l_i>1$}}
 \put(9,10){\vector(1,0){110}}
 \put(13,10){\circle*{1.2}}
 \put(58,10){\circle*{1.2}}
 \put(107,10){\circle*{1.2}}
 \put(8,12){\makebox(10,5){$y_i$}}
 \put(53,12){\makebox(10,5){$\zeta_i$}}
 \put(102,12){\makebox(10,5){$x$}}
 \put(8,3){\makebox(10,5){$0$}}
 \put(43,3){\makebox(30,5){$1/l_i<1$}}
 \put(102,3){\makebox(10,5){$1$}}
\end{picture}
\caption{Two independent reverse rays, shown schematically with their
duration parameters. The modified center endpoint $\widehat x$ and
the true outer endpoint $\zeta_i$ are strictly before cut. The full
endpoints $\zeta$ and $x$ may still be conjugate at this stage of the
argument. The two rows are not drawn to a common scale.}
\label{fig:reverse-rays}
\end{figure}

\subsection{The outward endpoint and conjugacy}

\begin{lemma}[Conjugacy is excluded by an active increasing modification]
\label{jet:outer-conjugacy}
Fix $b>0$ and an active outer pair
\[
 g_b(x)+c(x,y)+\varphi_o(v(y))=0.
\]
Choose any minimizing reverse log $r$ from $y$ to $x$. Suppose
\[
 l=\varphi_o'(v(y))>1,\qquad \varphi_o''(v(y))<0.
\]
Then the chosen modified ray is nonconjugate at its endpoint.
The true log $r/l$ is before cut.
\end{lemma}

\begin{proof}
The active inequality reads
$\varphi_o(v(h))\ge-c(x,h)-g_b(x)$.
For $1/l<s<1$, put $\xi_s=\exp_y(sr)$ and form the smooth upper
support
\[
 C_s(h)=c(h,\xi_s)/s+c(\xi_s,x)/(1-s).
\]
It has value $c(y,x)$ and gradient $-r$ at $y$. Since
$\varphi_o$ is increasing, $\varphi_o^{-1}(-C_s-g_b(x))$
is a smooth lower test for $v$ with gradient $r/l$.
That log is strictly before cut because it shortens a minimizing
ray. It also gives a global contact by
Theorem~\ref{geo:support}, but its before-cut property does not
depend on a density equation.

The scaled true relative Hessian of this test is
\begin{equation}\label{jet:outer-truncated}
 W_s=lA_y(r/l)-s^{-1}A_y(sr)
       -\frac{\varphi_o''(v(y))}{l^2}rr^*.
\end{equation}
Fix $s_0\in(1/l,1)$. For $s\ge s_0$, divided-action monotonicity
gives a positive definite lower bound for the difference of the
first two terms, independent of $s$. The last term is positive
semidefinite. If the endpoint at duration one were conjugate,
the index-form characterization of conjugacy, equivalently
Lemma~\ref{spec:radial}, would force an eigenvalue of
$s^{-1}A_y(sr)$ to tend to $-\infty$ as $s\uparrow1$.
Thus $\det W_s\to\infty$, while $W_s$ retains a positive
definite baseline. Lemma~\ref{jet:viscosity}, at the fixed true
log $r/l$, instead gives
\[
 \det W_s\le l^n(\Lambda/\lambda)/\sigma_y(r/l)<\infty.
\]
This is a contradiction.
\end{proof}

\begin{lemma}[Undoing the outer modification]
\label{jet:outer-test}
At a pair as in Lemma~\ref{jet:outer-conjugacy}, let
$p=\log_xy$ be the corresponding forward log,
$B=(d\exp_x)_p$, and let $H_o\ge0$ be a lower Taylor Hessian,
with zero linear part, of
$d\mapsto w_o(\exp_x(p+d))+\tfrac12|p+d|^2$.
Define
\begin{align}
 V_o={}&H_o+B^*\{lA_y(r/l)-A_y(r)\}B
             -\frac{\varphi_o''(v(y))}{l^2}pp^* .
 \label{jet:outer-form}
\end{align}
Then $V_o>0$ and
\begin{equation}\label{jet:outer-det}
 \det V_o\le
 l^n(\Lambda/\lambda)\,
       \frac{\sigma_x(p)^2}{\sigma_y(r/l)}.
\end{equation}
The middle term in \eqref{jet:outer-form} is positive definite.
On a compact family of active pairs where $b>0$ is bounded away
from zero, $\varphi_o'>1$ and $\varphi_o''<0$ hold with fixed
strict margins, this term has a uniform positive definite
baseline and the chosen exponential branches are uniformly
nonsingular.
\end{lemma}

\begin{proof}
Lemma~\ref{jet:outer-conjugacy} makes the endpoint coordinate map
locally invertible. Its smooth cost branch has endpoint Hessian
$A_y(r)$, and the zero linear coefficient fixes the gradient of
the induced lower test of $w_o$ to be $r$. Inverting
$\varphi_o$ and using the chain rule shows that
\eqref{jet:outer-form} is $l$ times the true relative lower-test
Hessian, pulled back by $B$. Gauss's Lemma gives the rank-one
term displayed there.

The divided-action difference between durations $1/l$ and $1$
is strictly positive definite, so $V_o>0$. Decrease the lower
Taylor Hessian by $\eps\Id$ to obtain a genuine lower test,
apply Lemma~\ref{jet:viscosity}, and let $\eps\downarrow0$.
Taking the determinant of the scaling and pullback gives
\eqref{jet:outer-det}.

For the uniformity assertion, include the minimizing log among
the parameters. The family is compact and active contact is a
closed condition. Lemma~\ref{jet:outer-conjugacy} excludes a
conjugate limit anywhere in that family; radial shortening
excludes true cut limits. The smooth branch matrices and the
positive divided-action difference therefore have the claimed
uniform bounds by compactness. Multiple ordinary branches
cause no difficulty: each specified nonconjugate minimizing
branch is a smooth upper branch and has been included among
the parameters.
\end{proof}

\subsection{The order of the limiting argument}

We now select points at which the preceding lower tests and the original
Jacobian equation can be combined. The following elementary fact will
justify differentiation through a short-time minimizer once its endpoint
is an Alexandrov point: neighboring endpoints need not be differentiable
points of the datum.

\begin{lemma}[Subgradients at an Alexandrov point]\label{jet:subgradient-differential}
If $f$ is semiconvex near $0$ and has the Alexandrov expansion
\[
 f(h)=f(0)+p\cdot h+\tfrac12 A(h,h)+o(|h|^2),
\]
then every nearby subgradient satisfies
\[
 \zeta=p+Ah+o(|h|),\qquad \zeta\in\partial f(h),
\]
uniformly over the choice of $\zeta$.
\end{lemma}

\begin{proof}
Adding a fixed quadratic reduces the assertion to a
convex function, and subtracting it later restores the displayed
formula. In the convex case write the remainder as $R(h)$, where
$|R(h)|\le\omega(|h|)|h|^2$ for a nondecreasing modulus
$\omega(r)\to0$. For $r=|h|>0$, a unit vector $e$, and
$0<\alpha\le1$, the subgradient inequalities at
$h\pm\alpha r e$ give
\[
 |(\zeta-p-Ah)\cdot e|
 \le \tfrac12\|A\|\alpha r
        +5\omega(2r)r/\alpha.
\]
For small $r$, choose
$\alpha=\sqrt{\max\{\omega(2r),r^2\}}$ and take the supremum
over unit $e$. This proves the uniform assertion.
\end{proof}

The selection has two stages. At each fixed prefix time we first pass
from nearby twice differentiable points to an exact maximum, where the
parameter inequality holds for every active representation. After
selecting an outer endpoint in that limit, we return to the nearby
points to obtain the Hessian bound and the original center Jacobian.
This order preserves the parameter inequality without assuming that a
prescribed active vertex persists under perturbation.

\begin{proposition}[The limiting matrix configuration]
\label{jet:determinants}
Fix the center template and the constants of
Lemmas~\ref{sc:excess}, \ref{sc:maximum}, \ref{sc:parameter}, and
\ref{sc:outward}. Choose $K$, then an outer template with curvature
less than $-K$ in the prescribed capped region. For all sufficiently
small scales in the sequence of Lemma~\ref{sc:plateau}, there are
a parameter $b\in[b_-,b_+]$, a pole $x$, center and active logs
$p,p_i\in G_x$, weights $m_i>0$ summing to one, and symmetric
matrices $H,L,V,V_i$ on $T_xM$, with the following properties.
For one selected index $i$,
\begin{gather}
 p=\sum_jm_jp_j,\qquad
 E=\Span\{p_j-p\}=\R e,\qquad |e|=1,\quad
 q=p\cdot e>0,\quad q^2\ge cD, \label{jet:limiting-geometry}\\
 m_i\ge m_*>0,\qquad p_i=p+d_ie,\quad d_i>0, \label{jet:limiting-outward}\\
 L\ge0,\qquad LE=0,\qquad H\le L,\qquad
 H\ge V+c\frac bD\Id-C\frac b{D^2}pp^*,\qquad V>0,
 \label{jet:limiting-center}\\
 V_i\ge m_iL+cK\frac b{D^2}p_ip_i^*,\qquad V_i>0.
 \label{jet:limiting-outer}
\end{gather}
The center modified log and the selected outer modified log are
nonconjugate. Let $r$ and $r_i$ denote the reverse logs of these
specified rays, let $y=\exp_xp$, $y_i=\exp_xp_i$, and set
$l=\varphi_c'(v(y))$, $l_i=\varphi_o'(v(y_i))$. The chosen true
limiting branches are also nonconjugate, and
\begin{equation}\label{jet:limiting-determinants}
 \det V\ge c\,\frac{\sigma_x(p)^2}{\sigma_y(r/l)},
 \qquad
 \det V_i\le C\,\frac{\sigma_x(p_i)^2}{\sigma_{y_i}(r_i/l_i)}.
\end{equation}
Here $c,C,m_*$ are independent of $K$ and the outer template.
The level and norm-increment conclusions of
Lemma~\ref{sc:outward} hold as well.
\end{proposition}

\begin{proof}
\emph{Step 1: fixed scales and exact maxima.}
Throughout the proof the original potential $v$, the scales, and
both templates are fixed. Constants used only to pass limits may
depend on all of them. Choose $t_k\uparrow1$ and interior maxima
$(z_k^0,b_k)$ supplied by Lemma~\ref{sc:maximum}, for
\[
 F_k(z)=Q_{t_k}g_{b_k}(z),\qquad
 G_k(z)=Q_{1-t_k}w_{c,b_k}(z).
\]
The functions are $C^{1,1}$ in space:
Proposition~\ref{geo:intermediate} applies to the first, and
Lemma~\ref{geo:short-time} to the second. At an exact maximum the
velocities align. If $x_k^0$ and $y_k^0$ are its two minimizers
and $p_k^0=t_k^{-1}\log_{x_k^0}z_k^0$, then
$y_k^0=\exp_{x_k^0}p_k^0$ along the concatenated minimizing ray.
Every finite positive active representation of $p_k^0$ satisfies
the parameter inequality of Lemma~\ref{sc:parameter}. The level
relations and common section bounds of Lemma~\ref{sc:levels}
hold for these representations.

\emph{Step 2: null sets are removed before Jensen sampling.}
Let $N$ be the null set for the original potential in
\eqref{jet:ae-jacobian}. For this fixed $k$, the final minimizer
map $Y_k(z)$ is locally Lipschitz and hence differentiable
almost everywhere \cite[Chapter~2, Theorem~1.4]{Simon2018}.
In each coordinate chart the
locally Lipschitz area formula
\cite[Chapter~2, Theorem~3.3]{Simon2018} implies that
\begin{equation}\label{jet:null-preimage}
 Z_k=\{z:DY_k(z)\text{ exists},\ \det DY_k(z)\ne0,\
                     Y_k(z)\in N\}
\end{equation}
is a null set. Indeed the integral of $|\det DY_k|$ on
$Y_k^{-1}(N)$ is zero. This argument does not require us to know
in advance where $DY_k$ is nonsingular.

Apply the semiconvex maximum Lemma
\cite[Appendix, Lemmas~A.3--A.4]{CIL1992} to $F_k+G_k$ at
$z_k^0$, using a small fourth-order strictness penalty if needed.
Its positive-measure conclusion allows us to choose a sequence
$z_{k,\nu}\to z_k^0$ off $Z_k$ and off all the null sets where
$F_k,G_k$, the prefix pole map, or $Y_k$ lack the derivatives
used below. We have
\begin{equation}\label{jet:Jensen-sequence}
 \grad(F_k+G_k)(z_{k,\nu})\longrightarrow0,\qquad
 D^2(F_k+G_k)(z_{k,\nu})\le \eta_{k,\nu}\Id,\quad
 \eta_{k,\nu}\longrightarrow0
\end{equation}
in the fixed chart.
At every such point choose a representation with at most $n+1$
active logs, padding by zero weights only for compactness.
Lemma~\ref{jet:transfer} applies at each point. Its matrices
$L_{k,\nu}$ are positive semidefinite and annihilate the active
differences, exactly; its endpoint lower jets are $m_{i,k,\nu}
L_{k,\nu}$. The matrices $H_{k,\nu}$ satisfy
\eqref{jet:Jensen-error}.

\emph{Step 3: take the first, fixed-prefix limits.}
For each fixed $k$, the local $C^{1,1}$ bounds are finite.
After taking a subsequence in $\nu$, all configuration parameters,
weights, and the two matrices $H_{k,\nu},L_{k,\nu}$ converge.
Denote the limits by $p_k,p_{i,k},m_{i,k},H_k,L_k$ and the
corresponding base points. Discard zero limiting weights.
The limiting endpoints are active at the exact maximum, and
their weighted sum is $p_k$. In particular the \emph{universal}
parameter inequality applies to this limiting representation;
we have not asserted forward persistence of any prescribed
active vertex. We retain
\begin{equation}\label{jet:fixed-prefix-limits}
 L_k\ge0,\qquad L_kE_k=0,\qquad H_k\le L_k,
 \quad E_k=\Span\{p_{i,k}-p_k:m_{i,k}>0\}.
\end{equation}
The quantitative proof of Lemma~\ref{jet:semibound} has an error
bounded by a fixed-scale constant times
$1-t_k+|\grad(F_k+G_k)(z_{k,\nu})|$. Passing $\nu\to\infty$
therefore gives
\[
 H_k\ge c\frac{b_k}{D}\Id
           -C\frac{b_k}{D^2}p_kp_k^*
           -C_{\mathrm{fixed}}(1-t_k)\Id.
\]
Since $b_k\ge b_->0$, for sufficiently large $k$ this implies,
after reducing $c$,
\begin{equation}\label{jet:absorbed-semibound}
 H_k\ge c\frac{b_k}{D}\Id
           -C\frac{b_k}{D^2}p_kp_k^*.
\end{equation}

If $\dim E_k\ge2$, a nonzero vector in $E_k\cap p_k^\perp$
contradicts \eqref{jet:fixed-prefix-limits} and
\eqref{jet:absorbed-semibound}. If $E_k=0$, all active logs in
the representation equal $p_k$ and all endpoints equal the
center. But $B_c\le B_o$ at every level, contradicting the
strict parameter inequality. Thus $E_k=\R e_k$. Orient $e_k$
so $q_k=p_k\cdot e_k>0$; the same two inequalities give
$q_k^2\ge cD$. Lemma~\ref{sc:outward} now supplies an index
with $m_{i,k}\ge m_*$, a strict outward displacement, and a
bounded template value. We choose an index carrying a fixed
share of the outward moment, as permitted by that Lemma, so
its displacement has a positive lower bound at the fixed
scales. Relabel it as index $1$.

\emph{Step 4: fix the outer branch before bounding any center Hessian.}
Pass to a subsequence in $k$ so that the selected configurations
and their weights converge. Their centers and selected
endpoints remain in the compact level ranges of
Lemma~\ref{sc:outward}. The selected limiting weight is at
least $m_*$, its displacement is strictly outward, and its
outer template value is below the prescribed cap. By imposing
the template conditions under a slightly larger fixed cap,
the limiting slope and curvature conditions have strict
margins:
$B_o'>0$ and $B_o''<-K$. The limiting selected pair is active.
Lemma~\ref{jet:outer-conjugacy} excludes conjugacy of its
modified branch, using only this activity and the original
contact-volume bound. Its shortened true log is before cut.
Lemma~\ref{jet:outer-test} consequently supplies a neighborhood
of this configuration with constants $\delta>0$ and $M<\infty$
such that its divided-action gain is at least $\delta\Id$
and its determinant upper bound is at most $M$.
These constants may depend on the fixed scales and templates.

For every remaining $k$, return to its already chosen
Jensen subsequence and choose $\nu=\nu(k)$ sufficiently large.
We require the configuration and matrices to be within $1/k$
of their fixed-prefix limits, the selected weight to be at
least $m_*/2$, and the selected endpoint to lie in that
neighborhood. We also require both errors in
\eqref{jet:Jensen-sequence}, including their transformed
versions in \eqref{jet:Jensen-error}, to be at most
$(1-t_k)^3/k$. All these requirements are compatible because
the limit in $\nu$ is taken with $k$ fixed.
For the rest of the proof, a subscript $k$ denotes this
diagonal sequence of actual Jensen points.

\emph{Step 5: the outer test bounds the exact pole matrix.}
At each diagonal point, Lemma~\ref{jet:transfer} gives the
actual lower Taylor Hessian
$H_{o,k}=m_{1,k}L_k\ge0$ at the selected endpoint. Use
Lemma~\ref{jet:outer-test} to undo the modification and call
the resulting form $V_{1,k}$. The positive outer curvature
term and the uniform divided-action baseline give
\begin{equation}\label{jet:bootstrap}
 V_{1,k}\ge m_{1,k}L_k+\delta\Id,\qquad
 \det V_{1,k}\le M.
\end{equation}
Every eigenvalue of $V_{1,k}$ is at least $\delta$, so its
largest eigenvalue is at most $M\delta^{1-n}$. Since
$m_{1,k}\ge m_*/2$, this bounds $L_k$ from above.
The comparison $H_k\le L_k+o(1)\Id$ and
Lemma~\ref{jet:semibound} now bound $H_k$ from both sides.
There has still been no use of a center density equation.

The center extension gives uniformly nonsingular
fixed-scale prefix-log coordinate maps, as in the proof of
Lemma~\ref{jet:semibound}. Their derivatives and second
derivatives are bounded; the gradients of $G_k$ are bounded
by the Lipschitz bound of the datum. Thus the coordinate
change in \eqref{jet:H-definition} converts the bound on
$H_k$ into
\begin{equation}\label{jet:G-bound}
 \|\Hess G_k(z_k)\|\le C_{\mathrm{fixed}}.
\end{equation}

\emph{Step 6: the chosen centers are good original endpoints.}
Write $\tau_k=1-t_k$, $y_k=Y_k(z_k)$, and
$K_k(z,y)=c(z,y)/\tau_k$. The envelope identity
$\grad G_k=(K_k)_z(z,Y_k(z))$ and differentiability of $Y_k$
give
\[
 \Hess G_k=(K_k)_{zz}+(K_k)_{zy}DY_k.
\]
In parallel orthonormal frames, the distance between $z_k$
and $y_k$ is $O(\tau_k)$, and the smooth short-cost expansions
are
\[
 (K_k)_{zz}=\tau_k^{-1}\Id+O(\tau_k),\quad
 (K_k)_{yy}=\tau_k^{-1}\Id+O(\tau_k),\quad
 (K_k)_{zy}=-\tau_k^{-1}\Pi+O(\tau_k),
\]
where $\Pi$ is the appropriate parallel isometry. Solving
the preceding identity, using \eqref{jet:G-bound}, gives
\begin{equation}\label{jet:Y-identity}
 DY_k=\Pi^{-1}\{\Id-\tau_k\Hess G_k\}+O(\tau_k^2).
\end{equation}
It is nonsingular for large $k$. The Jensen points were
chosen off the set \eqref{jet:null-preimage}; hence
$y_k\notin N$. This is the promised justification for
using the original Alexandrov and Jacobian statements at
the centers.

In the following formulas the scalar modifications are evaluated
with $b=b_k$. At these points $w_c=\varphi_c\circ v$ has an Alexandrov
Hessian. The minimizing property supplies a proximal subgradient
\[
 \xi(z)=-(K_k)_y(z,Y_k(z))\in\partial w_c(Y_k(z))
\]
at every nearby point, even if $w_c$ is not differentiable there.
Apply Lemma~\ref{jet:subgradient-differential} at the Alexandrov point
$y_k=Y_k(z_k)$. Since $Y_k$ is Lipschitz and differentiable at
$z_k$, it gives, in normal coordinates at $z_k$ and $y_k$,
\[
 \xi(z_k+h)=\grad w_c(y_k)
       +\Hess w_c(y_k)DY_k(z_k)h+o(|h|).
\]
Expanding the smooth function $(K_k)_y$ in its two variables
therefore gives
\[
 \{\Hess w_c(y_k)+(K_k)_{yy}\}DY_k+(K_k)_{yz}=0.
\]
Thus stationary differentiation uses the subgradient expansion even
when neighboring endpoints are nonsmooth.
Solving it with \eqref{jet:Y-identity} gives, in parallel
frames,
\begin{equation}\label{jet:undo-short}
 \Hess w_c(y_k)=\Hess G_k(z_k)+O(\tau_k).
\end{equation}
In particular these endpoint Hessians are uniformly bounded.
The error statement follows directly from the displayed
short-cost expansions and the bounded inverse in
\eqref{jet:Y-identity}; it does not assume a bound for
$\Hess w_c$ in advance.

\emph{Step 7: insert the true center form and take matrix limits.}
Set
\[
 r_k=\tau_k^{-1}\log_{y_k}z_k,\quad
 l_k=\varphi_c'(v(y_k)),\quad
 \widehat x_k=\exp_{y_k}r_k,\quad
 \widehat p_k=\log_{\widehat x_k}y_k.
\]
The true center endpoint is
$\zeta_k=\exp_{y_k}(r_k/l_k)$, and it is before cut because
$y_k\notin N$. Identify $T_{x_k}M$ with
$T_{\widehat x_k}M$ by parallel transport over their short
joining segment, and let $B_k:T_{x_k}M\to T_{y_k}M$ be the
exponential differential at $\widehat p_k$ after that
identification. Define
\begin{equation}\label{jet:V-sequence}
 V_k=l_k B_k^*
       \{\Hess v(y_k)+A_{y_k}(r_k/l_k)\}B_k>0.
\end{equation}
The exact chain rule for the scalar modification gives
\begin{align}
 B_k^*\{\Hess w_c(y_k)+A_{y_k}(r_k)\}B_k
 ={}& V_k
 +B_k^*\{A_{y_k}(r_k)-l_kA_{y_k}(r_k/l_k)\}B_k
 \nonumber\\
 &+\frac{\varphi_c''(v(y_k))}{l_k^2}
      \widehat p_k\widehat p_k^* .
 \label{jet:center-exact}
\end{align}
All vectors on the last line are understood after the
same isometric identification.
Lemma~\ref{spec:pullback} bounds the middle term below by
$c(1-l_k)\Id$. Equation~\eqref{jet:undo-short} and the
vanishing pole mismatch show that the left side differs
from $H_k$ by $o(1)\Id$. The verification is the same smooth
coordinate comparison as in
Lemma~\ref{jet:semibound}, now using the bounded
unregularized Hessians in place of the smooth support
Hessians. This proves
\[
 H_k\ge V_k+c\frac{b_k}{D}\Id
       -C\frac{b_k}{D^2}p_kp_k^*-o(1)\Id.
\]
It also bounds $V_k$ from above, because $V_k>0$ and $H_k$
is bounded.

Taking the determinant in \eqref{jet:V-sequence} and using
\eqref{jet:ae-jacobian} gives the exact factorization
\begin{equation}\label{jet:V-jacobian}
 \det V_k
 =l_k^n\sigma_{\widehat x_k}(\widehat p_k)^2
       \frac{|\det D(\exp_{\cdot}\grad v)(y_k)|}
            {\sigma_{y_k}(r_k/l_k)}
 \ge c\,\frac{\sigma_{\widehat x_k}(\widehat p_k)^2}
                 {\sigma_{y_k}(r_k/l_k)}.
\end{equation}
The squared factor is the determinant of the Hessian
pullback. All scalar derivatives lie between $1/2$ and $2$,
so the constants in this density inequality are independent
of $K$.

There is no assertion that the limiting endpoint is an
Alexandrov point of $v$. We take limits of matrices at the good
endpoints just selected. The center modified rays have a uniform
extension at the fixed scales: $1-l_k\ge c b_-/D>0$, and the
ray through $\widehat x_k$ continues to the true endpoint through
duration $1/l_k$. Compactness of this family of proper minimizing
subsegments therefore gives
\[
 \sigma_{\widehat x_k}(\widehat p_k)\ge c_{\mathrm{fixed}}>0.
\]
Since $0<V_k\le C_{\mathrm{fixed}}\Id$, equation
\eqref{jet:V-jacobian} now implies
\[
 \sigma_{y_k}(r_k/l_k)
 \ge \frac{c\,\sigma_{\widehat x_k}(\widehat p_k)^2}
             {\det V_k}
 \ge c'_{\mathrm{fixed}}>0.
\]
Continuity of the exponential differential excludes a conjugate
limit of the true center rays. An ordinary cut limit may remain;
the specified nonsingular branch suffices.

Finally take a subsequence so that $b_k,H_k,L_k,V_k,V_{1,k}$
and all configuration parameters converge, in a smooth
orthonormal trivialization. Discard any additional zero
limiting weights. The selected weight remains at least $m_*$,
and its outward displacement retains the positive fixed-scale
lower bound from Step~3. The barycenter identity then gives
another surviving vertex, so the surviving difference span
is still the same nonzero line. The selected modified and true
outer branches were already uniformly nonsingular in
Step~4; the center modified branch has its fixed-scale
extension. Equation~\eqref{jet:V-jacobian} gives a positive
lower determinant in the limit, so $V>0$. Equations
\eqref{jet:L-properties}, \eqref{jet:Jensen-error}, and
\eqref{jet:center-exact} give
\eqref{jet:limiting-center}.

At each selected outer endpoint in the diagonal sequence, formula
\eqref{jet:outer-form}, with $H_{o,k}=m_{1,k}L_k$, has a positive
divided-action term. Its scalar curvature coefficient satisfies
\[
 -\frac{\varphi_o''(v(y_{1,k}))}{l_{1,k}^2}
 =-\frac{b_kB_o''((v(y_{1,k})-a)/D)}{D^2l_{1,k}^2}
 \ge cK b_k/D^2.
 \]
Passing this inequality and \eqref{jet:outer-det} along
$V_{1,k}\to V_i$ gives \eqref{jet:limiting-outer} and the other
inequality in \eqref{jet:limiting-determinants}. No lower Taylor
jet is asserted at the limiting endpoint. The weights, directions,
and section-level estimates pass through the same chosen
subsequence and give
\eqref{jet:limiting-geometry}--\eqref{jet:limiting-outward}.
Their constants and the cap were fixed before $K$.
Only $\delta,M$, the rate of the prefix limit, and the
Jensen radii depended on the fixed outer template.
\end{proof}

\subsection{The determinant gain}

\begin{lemma}[Oblique projection and the missing direction]
\label{jet:matrix}
Suppose the matrices and vectors satisfy
\eqref{jet:limiting-geometry}--\eqref{jet:limiting-outer}.
Then
\begin{equation}\label{jet:matrix-gain}
 \det V_i\ge cK\det V,
\end{equation}
where $c>0$ depends on the fixed center-template constants,
the dimension, and $m_*$, and is independent of $K$.
\end{lemma}

\begin{proof}
Put $\beta=Cb/D^2$, increasing $C$ if necessary, and
$\gamma=cKb/D^2$. Define
\[
 P=\Id-e p^*/q.
\]
This is the projection onto $p^\perp$ along $e$.
Since $LE=0$, $P^*LP=L$. On $p^\perp$ the center inequality
gives $H\ge V$, and $L\ge H$. Thus
\begin{equation}\label{jet:projected-matrix}
 V_i\ge m_iP^*VP+\gamma p_ip_i^*.
\end{equation}
Also $H(e,e)\le L(e,e)=0$ implies
$V(e,e)\le\beta q^2$.

Let $Q$ have orthonormal columns spanning $p^\perp$, and
use the nonsingular basis matrix $[Q,e]$. In this basis write
\[
 [Q,e]^*V[Q,e]
 =\begin{pmatrix} A&z\\z^*&h\end{pmatrix},\qquad
 S=h-z^*A^{-1}z>0.
\]
Then $S\le h=V(e,e)\le\beta q^2$.
The determinant of the right side of
\eqref{jet:projected-matrix}, in the same basis, is
$m_i^{n-1}\det A\,\gamma(p_i\cdot e)^2$.
Taking the ratio cancels the determinant of the
nonorthogonal change of basis and gives
\[
 \frac{\det V_i}{\det V}
 \ge m_i^{n-1}\frac{\gamma(p_i\cdot e)^2}{S}
 \ge m_i^{n-1}\frac{\gamma}{\beta}
        \frac{(p_i\cdot e)^2}{q^2}.
\]
Since $p_i=p+d_ie$ with $d_i>0$ and $q>0$, we have
$p_i\cdot e\ge q$. Now use $m_i\ge m_*$ and
$\gamma/\beta\ge cK$.
\end{proof}

\begin{proposition}[Contradiction to a non-power oscillation]
\label{jet:contradiction}
The section-oscillation function of Lemma~\ref{sc:plateau}
has a positive power upper bound at zero.
\end{proposition}

\begin{proof}
Suppose otherwise, and take the sequence of scales supplied
by that Lemma. Fix all center and weight constants, then
choose $K$ and its outer template, and then take sufficiently
small scales to apply Proposition~\ref{jet:determinants}.
Lemma~\ref{jet:matrix} and
\eqref{jet:limiting-determinants} imply
\[
 cK\,\frac{\sigma_x(p)^2}{\sigma_y(r/l)}
 \le \det V_i
 \le C\,\frac{\sigma_x(p_i)^2}{\sigma_{y_i}(r_i/l_i)}.
\]
The radial comparison of Corollary~\ref{spec:radial-determinants} gives
\[
 \sigma_y(r/l)\le C\sigma_x(p),\qquad
 \sigma_{y_i}(r_i/l_i)\ge c\sigma_x(p_i).
\]
For the selected outward endpoint,
\eqref{jet:limiting-geometry}--\eqref{jet:limiting-outward} give
$m_i\ge m_*$ and $q^2\ge cD$. Its norm increment satisfies
$|p_i|^2-|p|^2\le3D$ by Lemma~\ref{sc:outward}.
Corollary~\ref{spec:outward} therefore applies with
$\mu=m_*$, $\kappa=c$, and $L_0=3$, and gives
$\sigma_x(p)\ge c'\sigma_x(p_i)$. These parameters were fixed
before $K$. Hence
\[
 cK\sigma_x(p)\le C\sigma_x(p_i)\le C'\sigma_x(p).
\]
The center modified branch is nonconjugate, so
$\sigma_x(p)>0$. This bounds $K$ by a constant fixed before
it was chosen. Choosing $K$ larger gives a contradiction.
\end{proof}
